\documentclass[french,a4paper]{article}

\usepackage[utf8]{inputenc}
\usepackage[T1]{fontenc}
\usepackage{lmodern}
\usepackage{geometry}
\usepackage{graphicx}
\usepackage{babel}
\usepackage{amsmath}
\usepackage{amsthm}
\usepackage{amssymb}
\usepackage{hyperref}
\usepackage{stmaryrd}
\usepackage{enumitem}
\usepackage{tcolorbox}
\usepackage{dsfont}
\usepackage{multirow}
\usepackage{etoc}
\usepackage{titlesec}
\usepackage{esint}
\usepackage{cancel}
\usepackage{mathdots}

\setlist[itemize]{label=$\bullet$}


\renewcommand{\P}{\mathbb{P}}
\newcommand{\N}{\mathbb{N}}
\newcommand{\Z}{\mathbb{Z}}
\newcommand{\Q}{\mathbb{Q}}
\newcommand{\R}{\mathbb{R}}
\newcommand{\C}{\mathbb{C}}
\newcommand{\K}{\mathbb{K}}
\renewcommand{\L}{\mathbb{L}}
\newcommand{\U}{\mathbb{U}}
\newcommand{\st}{^*}
\newcommand{\tq}{\middle|}
\newcommand{\inv}{^{-1}}
\newcommand{\E}{\mathbb{E}}
\newcommand{\F}{\mathbb{F}}
\newcommand{\V}{\mathbb{V}}
\renewcommand{\ker}{\text{Ker}}
\newcommand{\im}{\text{Im}}
\newcommand{\card}{\text{Card}}
\newcommand{\Tr}{\text{Tr}}

\theoremstyle{plain}
\newtheorem{thm}{Théorème}
\newtheorem*{ind}{Indication}

\theoremstyle{definition}
\newtheorem{dfn}{Definition}

\theoremstyle{remark}
\newtheorem*{rmq}{Remarque}
\newtheorem*{app}{Application}

\newtcolorbox[auto counter]{ex}[2][]{colback=white, colframe=green!50!white, coltitle=black, fonttitle=\bfseries, title={Exercice~\thetcbcounter~: #2}, after title={\hfill #1}}

\title{TD Matrices, déterminants (révisions)}

\begin{document}

\maketitle

Dans toute la feuille d'exercices, $\K$ désigne un sur-corps de $\Q$.

\section{Espaces vectoriels de matrices}

\begin{ex}[Moulin.Mat.01]{}
	Soit $M\in\mathcal{M}_{n,p}(\K)$. Déterminer le sous-espace vectoriel de $\mathcal{M}_{n,p}(\K)$ engendré par les matrices équivalentes à $M$.
\end{ex}

\begin{ex}[Moulin.Mat.02]{}
	Déterminer le sous-espace vectoriel de $\mathcal{M}_n(\K)$ engendré par les matrices de projecteurs.
\end{ex}
\begin{proof}
	Les $E_{i,i}$ sont des matrices de projecteur. De plus, si $i\ne j$, $E_{i,i}+E_{i,j}$ est aussi une matrice de projecteur. Donc les $E_{i,j}$ pour $i\ne j$ sont aussi dedans par différence. Le sous-espace vectoriel cherché est donc $\mathcal{M}_n(\K)$.
\end{proof}

\begin{ex}[Moulin.Mat.03]{Matrices damiers}
	\begin{enumerate}
		\item Caractériser en termes d'endomorphismes les ensembles :
		$$D_0=\{(a_{i,j})\in\mathcal{M}_n(\K)\ : \ i+j\equiv 1\ [2]\implies a_{i,j}=0\}$$
		$$D_1=\{(a_{i,j})\in\mathcal{M}_n(\K)\ : \ i+j\equiv 0\ [2]\implies a_{i,j}=0\}$$
		\item Montrer que $D_0$ et $D_1$ sont des sous-espaces vectoriels de $\mathcal{M}_n(\K)$ supplémentaires.
		\item Que peut-on dire du produit d'un élément de $D_i$ par un élément de $D_j$ ?
		\item Montrer que l'inverse d'un élément inversible de $D_i$ est dans $D_i$.
	\end{enumerate}
\end{ex}

\begin{ex}[Moulin.Mat.04]{Matrices magiques}
	Soit $n\geqslant 3$. Pour $A=(a_{i,j})\in\mathcal{M}_n(\C)$, on note : $$l_i(A)=\sum_{j=1}^{n}a_{i,j}\ c_j(A)=\sum_{i=1}^{n}a_{i,j}\ d(A)=\sum_{k=1}^{n}a_{k,k}\ a(A)=\sum_{k=1}^{n}a_{k,n+1-k}$$
	\begin{enumerate}
		\item Montrer que l'ensemble $\mathcal{N}_0$ des matrices $A$ vérifiant : $$\forall (i,j)\in\llbracket1,n\rrbracket^2,\ l_i(A)=c_j(A)=0$$ est un sous-espace vectoriel de $\mathcal{M}_n(\C)$ de dimension $(n-1)^2$. On pourra rechercher son intersection avec un sous-espace vectoriel de dimension $2n-1$ bien choisi.
		\item $\star$ Montrer que l'ensemble $\mathcal{M}_0$ des matrices $A$ vérifiant \textbf{de plus} $d(A)=a(A)=0$ est un sous-espace vectoriel de $\mathcal{M}_n(\C)$ de dimension $(n-1)^2-2$.
		\item Quelle est la dimension du sous-espace vectoriel $\mathcal{M}$ des matrices $A$ vérifiant : $$\forall (i,j)\in\llbracket1,n\rrbracket^2,\ l_i(A)=c_j(A)=d(A)=a(A)\ ?$$
	\end{enumerate}
\end{ex}

\begin{ex}[Moulin.Mat.05]{Quaternions}
	Montrer que $$\mathbb{H}=\left\{\begin{pmatrix}
		a&-\overline{b}\\
		b&\overline{a}
	\end{pmatrix}\bigg| (a,b)\in\C^2\right\}$$ est un sous-espace vectoriel réel de $\mathcal{M}_2(\C)$ mais pas un sous-espace vectoriel complexe. Montrer que $\mathbb{H}$ est une $\R$-algèbre dans laquelle tout élément non nul est inversible. Est-ce un corps ?
\end{ex}

\begin{ex}[Moulin.Mat.06]{$\Delta$ Idéaux bilatères de $\mathcal{M}_n(\K)$}
	Un idéal bilatère de $\mathcal{M}_n(\K)$ est un sous-espace vectoriel $\mathcal{I}$ de $\mathcal{M}_n(\K)$ tel que : $$\forall A\in\mathcal{I},\ \forall M\in\mathcal{M}_n(\K),\ AM\in\mathcal{I} \ \ \text{et} \ \ MA\in\mathcal{I}$$
	\begin{enumerate}
		\item Montrer qu'un idéal bilatère non nul contient une matrice de rang $1$, puis qu'il contient toutes les matrices de rang $1$.
		\item Déterminer l'ensemble des idéaux bilatères de $\mathcal{M}_n(\K)$.
		\item Montrer que tout endomorphisme d'algèbre de $\mathcal{M}_n(\K)$ est injectif.
	\end{enumerate}
\end{ex}
\begin{proof}
	\begin{enumerate}
		\item Il contient une matrice non nulle, et on se ramène au fait qu'il contient une matrice $J_r$. Ensuite, on peut faire des opérations élémentaires matriciellement pour se ramener à $J_1$. Puis cet idéal contient toute matrice de rang 1 en multipliant par des matrices inversibles à droite et à gauche. Donc il contient tous les $E_{i,i}$ et est alors égal à $\mathcal{M}_n(\K)$.
		\item Ce sont donc exactement $\{0\}$ et $\mathcal{M}_n(\K)$ tout entier.
		\item Le noyau d'un tel morphisme est un idéal bilatère. Or, ce noyau n'est pas égal à $\mathcal{M}_n(\K)$ puisque l'image de $I_n$ est $I_n$. Donc son noyau est $\{0\}$ et ce morphisme est bien injectif.
	\end{enumerate}
\end{proof}

\section{Calcul matriciel, représentation matricielle}

\begin{ex}[Moulin.Mat.07]{$\circ$}
	Résoudre dans $\mathcal{M}_2(\R)^2$ les systèmes :
	\begin{enumerate}
		\item $XY=\displaystyle\begin{pmatrix}
			0&1\\
			0&0
		\end{pmatrix}$ et $YX=\displaystyle\begin{pmatrix}
		0&0\\
		0&0
		\end{pmatrix}$ ;
		\item $XY=\displaystyle\begin{pmatrix}
			1&0\\
			0&0
		\end{pmatrix}$ et $YX=\displaystyle\begin{pmatrix}
			0&0\\
			0&0
		\end{pmatrix}$.
	\end{enumerate}
\end{ex}

\begin{ex}[Moulin.Mat.08]{$\Delta$}
	Calculer les puissances $A^k$ de la matrice $A=\displaystyle\left(\binom{j}{i}\right)_{0\leqslant i,j\leqslant n}$ pour $k\in\Z$.
\end{ex}
\begin{proof}
	Pour cela, on interprète cette matrice comme la matrice de l'endomorphisme $$\begin{array}{lrcl}
		T:&\K_n[X]&\longrightarrow&\K_n[X]\\
		&P&\longmapsto&P(X+1)
	\end{array}$$ dans la base canonique de $\K_n[X]$. Par conséquent, la puissance $k$ de $A$ correspond à $P\mapsto P(X+k)$, que l'on sait très facilement calculer.
\end{proof}

\begin{ex}[Moulin.Mat.09]{$\circ$}
	Soit $A\in\mathcal{M}_n(\K)$ une matrice triangulaire supérieure stricte. Calculer $A^n$ ? A quelle condition nécessaire et suffisante une matrice triangulaire est-elle nilpotente ?
\end{ex}

\begin{ex}[Moulin.Mat.10]{$\circ$}
	Soit $E$ un espace vectoriel de dimension finie $n$, ainsi que $F$ et $G$ des sous-espaces vectoriels de $E$ de dimensions respectives $p$ et $q$. Calculer la dimension de : $$\{u\in\mathcal{L}(E)\ :\ u(F)\subset G\}$$
\end{ex}

\begin{ex}[Moulin.Mat.11]{}
	Soit $u$ un endomorphisme de rang $r$ d'un espace vectoriel $E$ de dimension finie $n$. On pose : $$A=\{v\in\mathcal{L}(E)\ : \ u\circ v=0\}\ \ \text{et}\ \ B=\{v\in\mathcal{L}(E)\ : \ v\circ u=0\}$$ Déterminer les dimensions de $A$, $B$ et $C=A\cap B$.
\end{ex}

\begin{ex}[Moulin.Mat.12]{$\Delta$ Matrice antédiagonale}
	Soient $a_1,\ldots,a_n$ des complexes non nuls. Donner l'inverse de la matrice antédiagonale : $$A=\begin{pmatrix}
		(0)&&a_n \\
		&\iddots&\\
		a_1&&(0)
	\end{pmatrix}$$
\end{ex}
\begin{proof}
	Son inverse est la matrice antédiagonale $$B=\begin{pmatrix}
		(0)&&a_1\inv \\
		&\iddots&\\
		a_n\inv&&(0)
	\end{pmatrix}$$ On peut le vérifier par un calcul direct ou le voir en interprétant la matrice antédiagonale comme la matrice d'une inversion des vecteurs de la base, pondérée par les coefficients $a_i$.
\end{proof}

\begin{ex}[Moulin.Mat.13]{Matrice Attila}
	On note $J\in\mathcal{M}_n(\K)$ la matrice Attila dont tous les coefficients valent $1$.
	\begin{enumerate}
		\item Déterminer la sous-algèbre de $\mathcal{M}_n(\K)$ engendrée par $J$.
		\item Quels sont les éléments inversibles de cette sous-algèbre ? Quels sont les rangs des éléments non inversibles ?
		\item Calculer $(I+\alpha J)^p$ pour $p\in\N$ puis pour $p\in\Z$ lorsque cela a un sens.
	\end{enumerate}
\end{ex}

\begin{ex}[Moulin.Mat.14]{}
	Soit $A$, $B$ et $C$ trois matrices carrées de même taille non nulles.
	\begin{enumerate}
		\item Montrer que si $ABC=0$, alors deux des trois matrices $A$, $B$ et $C$ sont non inversibles.
		\item Donner une condition nécessaire et suffisante sur les rangs de $A$ et $C$ pour qu'il existe $B$ inversible telle que $ABC=0$.
	\end{enumerate}
\end{ex}

\begin{ex}[Moulin.Mat.15]{}
	Déterminer toutes les matrices $A$ inversibles de $\mathcal{M}_n(\R)$ telles que $A$ et $A\inv$ soient à coefficients positifs ou nuls.
\end{ex}

\begin{ex}[Moulin.Mat.16]{$\star$}
	Soit $A$ et $B$ dans $\mathcal{M}_n(\K)$. On suppose qu'il existe au moins $n+1$ scalaires $\lambda_0,\ldots,\lambda_n\in\K$ pour lesquels $A+\lambda B$ est nilpotente. Montrer que $A$ et $B$ sont nilpotentes.
\end{ex}
\begin{proof}
	Remarquons que les $A+\lambda_i B$ sont nécessairement d'indice de nilpotence inférieure à $n$. On considère $$P=(A+X.B)^n\in\mathcal{M}_n(\K(X))$$ On note $P_{i,j}$ le coefficient d'indices $(i,j)$ de $P$. Alors tous les $P_{i,j}$ sont de degré au plus $n$ et s'annulent en $n+1$ points distincts, donc tous les $P_{i,j}$ sont nuls. Donc $P$ est la matrice nulle. En particulier, en substituant $0$ à $X$, on en déduit que $A$ est nilpotente. En factorisant de force dans $P(\lambda)$ par $\lambda$ et puisqu'il y a une infinité de scalaires non nuls, on en que $\tilde{P}=(X.A+B)^n$ est aussi nul, donc que $B$ est nilpotente.
\end{proof}

\begin{ex}[Moulin.Mat.17]{}
	Soit $M=\displaystyle\begin{pmatrix}
		A&B\\
		0&D
	\end{pmatrix}$ avec $A$ et $D$ des matrices carrées. Montrer que $M$ est inversible si, et seulement si $A$ et $D$ le sont. Déterminer dans ce cas l'inverse de $M$.
\end{ex}

\begin{ex}[Moulin.Mat.18]{}
	Soit $A$ et $B$ dans $\mathcal{M}_n(\C)$ de même rang telles que $A^2B=A$. Montrer que $B^2A=B$.
\end{ex}
\begin{proof}
	Traduire les hypothèses en termes d'applications linéaires. En déduire que $\im(a)$ et $\ker(a)$ sont supplémentaires, puis raisonner matriciellement.
\end{proof}

\section{Rang, équivalence}

\begin{ex}[Moulin.Mat.19]{$\circ$}
	\begin{enumerate}
		\item Quel est le rang de la matrice $(i+j-1)_{1\leqslant i,j\leqslant n}$ ?
		\item Quel est le rang de la matrice $((\alpha+i+j)^2)_{1\leqslant i,j\leqslant n}$ ?
	\end{enumerate}
\end{ex}

\begin{ex}[Moulin.Mat.20]{}
	Pour $A=(a_{i,j})\in\mathcal{M}_{n,p}(\K)$ et $B\in\mathcal{M}_{m,q}(\K)$, on note : $$A\otimes B=\begin{pmatrix}
		a_{1,1}B&a_{1,2}B&\cdots&a_{1,n}B\\
		a_{2,1}B&a_{2,2}B&\cdots&a_{2,n}B\\
		\vdots&\vdots&&\vdots\\
		a_{n,1}B&a_{n,2}B&\cdots&a_{n,n}B
	\end{pmatrix}\in\mathcal{M}_{mn,pq}(\K)$$
	Déterminer le rang de $A\otimes B$.
\end{ex}

\begin{ex}[Moulin.Mat.21]{$\star$}
	Soit $A\in\mathcal{M}_{p,n}(\K)$ et $B\in\mathcal{M}_{m,q}(\K)$. Déterminer le rang de l'application de $\mathcal{M}_{n,m}(\K)$ dans $\mathcal{M}_{p,q}(\K)$ qui à $M$ associe $AMB$. Retrouver alors le résultat de l'exercice précédent.
\end{ex}

\begin{ex}[Moulin.Mat.22]{$\circ$}
	Pour $A\in\mathcal{M}_{n,p}(\C)$, déterminer le rang de $\overline{A}$.
\end{ex}

\begin{ex}[Moulin.Mat.23]{}
	Quelles sont les matrices $A\in\mathcal{M}_n(\K)$ telles que : $$\forall M\in\mathcal{M}_n(\K),\ \text{rg}(AM)=\text{rg}(M)\ ?$$
\end{ex}

\begin{ex}[Moulin.Mat.24]{}
	Quelles sont les matrices $A\in\mathcal{M}_n(\K)$ telles que : $$\forall B\in\mathcal{M}_n(\K),\ \text{rg}(AB)=\text{rg}(BA)\ ?$$
\end{ex}

\begin{ex}[Moulin.Mat.25]{$\star$}
	Trouver tous les sous-espaces vectoriels de dimension $n$ de $\mathcal{M}_n(\K)$ ne contenant que des matrices de rang inférieur ou égal à $1$.
\end{ex}

\begin{ex}[Moulin.Mat.26]{$\star$}
	Soit $X_1,\ldots,X_r$, $Y_1,\ldots,Y_r$ des éléments non nuls de $\K^n$. Montrer que la matrice $M=\displaystyle\sum_{i=1}^{r}X_iY_i^T$ est de rang inférieur ou égal à $r$ et qu'elle est de rang $r$ si, et seulement si, les deux familles $(X_1,\ldots,X_r)$ et $(Y_1,\ldots,Y_r)$ sont libres.
\end{ex}

\begin{ex}[Moulin.Mat.27]{}
	Soit $M=\displaystyle\begin{pmatrix}
		A&C\\
		0&B
	\end{pmatrix}$ une matrice triangulaire par blocs. Comparer le rang de $M$ à $\text{rg}(A)+\text{rg}(B)$. Peut-on préciser lorsque $A$ ou $B$ est inversible ?
\end{ex}

\begin{ex}[Moulin.Mat.28]{}
		Soit $M=\displaystyle\begin{pmatrix}
		I_n&A\\
		A&0
	\end{pmatrix}$ avec $A\in\mathcal{M}_n(\K)$. Montrer que $\text{rg}(M)=n+\text{rg}(A^2)$.
\end{ex}
Vérifier que le carré est bien sur $A$ et non pas sur son rang.

\begin{ex}[Moulin.Mat.29]{}
	Soit $A\in\mathcal{M}_{p,n}(\K)$ et $B\in\mathcal{M}_{n,p}(\K)$. Montrer : $$\text{rg}\begin{pmatrix}
		I_n&B\\
		A&I_p
	\end{pmatrix}=\text{rg}(I_p-AB)+n=\text{rg}(I_n-BA)+p$$
\end{ex}

\begin{ex}[Moulin.Mat.30]{$\star$ Une application combinatoire du rang}
	\begin{enumerate}
		\item Soit $M=(m_{i,j})\in\mathcal{M}_n(\R)$ et $r>0$. On suppose :
		\begin{itemize}
			\item $\forall (i,j)\in\llbracket1,n\rrbracket^2,\ i\ne j\implies m_{i,j}=r$ ;
			\item $\forall i\in\llbracket1,n\rrbracket,\ m_{i,i}\geqslant r$ ;
			\item $\card\{i\in\llbracket1,n\rrbracket\ : \ m_{i,i}=r\}\leqslant 1$.
		\end{itemize} Montrer que $M$ est inversible.
		\item Soit $m$, $n$ et $r$ trois entiers strictement positifs, $E=\{a_1,\ldots,a_n\}$ un ensemble à $m$ éléments et $(A_j)_{1\leqslant j\leqslant n}$ une famille de $n$ parties de $E$ deux à deux distinctes. Écrire la matrice $M=(\card(A_i\cap A_j))\in\mathcal{M}_n(\R)$ sous la forme $B^TB$ où $B$ est une matrice de $\mathcal{M}_{m,n}(\R)$ à préciser.
		\item En déduire que si : $$\forall (i,j)\in\llbracket1,n\rrbracket^2,\ i\ne j\implies \card(A_i\cap A_j)=r$$ alors $n\leqslant m$.
	\end{enumerate}
\end{ex}

\begin{ex}[Moulin.Mat.31]{}
	Soit $\K$ un sous-corps de $\L$.
	\begin{enumerate}
		\item Comparer le rang dans $\mathcal{M}_{n,p}(\K)$ et le rang dans $\mathcal{M}_{n,p}(\L)$ d'une matrice $M$ de $\mathcal{M}_{n,p}(\K)$.
		\item Soit $A\in\mathcal{M}_n(\K)$. On note : $$\mathcal{C}_\K=\{M\in\mathcal{M}_n(\K)\ : \ AM=MA\}\ \ \text{et}\ \ \mathcal{C}_\L=\{M\in\mathcal{M}_n(\L)\ : \ AM=MA\}$$ Comparer $\dim_\K(\mathcal{C}_\K)$ et $\dim_\L(\mathcal{C}_\L)$.
	\end{enumerate}
\end{ex}

\begin{ex}[Moulin.Mat.32]{}
	Soit $f:\mathcal{M}_n(\C)\to\C$ non constante telle que : $$\forall(A,B)\in\mathcal{M}_n(\C)^2,\ f(AB)=f(A)f(B)$$
	\begin{enumerate}
		\item Montrer que $A\in\mathcal{M}_n(\C)$ est inversible si, et seulement si, $f(A)\ne 0$. On pourra penser aux matrices nilpotentes.
		\item Trouver tous les $X\in\mathcal{M}_n(\C)$ tels que : $$\forall M\in\mathcal{M}_n(\C),\ f(M+X)=f(M)+f(X)$$
	\end{enumerate}
\end{ex}
\begin{proof}
	Déjà on remarque que $f(0)=f(0)^2$ donc $f(0)=0$ ou $f(0)=1$. De même pour $I_n$. Or, il est impossible que $f(0)=0$ sinon $f$ serait constante égale à $f(0)$ car $0$ est absorbant pour $\times$. De même, on a nécessairement $f(I_n)=1$.
	\begin{enumerate}
		\item Dans le sens direct, c'est immédiat car alors $f(A)f(A\inv)=1$ donc les deux sont non nuls. Dans le sens réciproque, on raisonne par contraposée. Si $A$ est non inversible, $A$ est équivalente à une matrice $J_r$ avec $r<n$. Mais cette $J_r$ est équivalente à une matrice $B_r$  dont la diagonale supérieure ne contient que $r$ 1 et des 0, qui est nilpotente. Par conséquent, $f(B_r)^n=0$. Donc $f(A)^n=0$ par commutativité de $\C$. Donc $f(A)=0$.
		\item Nécessairement, en écrivant $X=PJ_rQ$ et $M=P(I_n-J_r)Q$, on obtient par hypothèse $r=0$ ou $r=n$, donc $X+0$ ou $X\in\mathcal{GL}_n(\C)$. Réciproquement, $0$ convient, mais pour les matrices inversibles c'est pkus compliqué. $I_n$ ne convient pas lorsque $n\geqslant2$ (prendre $-E_{1,1}$) et si $n\geqslant 1$, on ne peut pas dire grand-chose.
	\end{enumerate}
\end{proof}
\section{Similitude}

\begin{ex}[Moulin.Mat.33]{$\circ$}
	Soit $u$ un endomorphisme d'un espace vectoriel $E$ de dimension finie. Montrer que les conditions suivantes sont équivalentes :
	\begin{enumerate}[label=(\roman*)]
		\item $E=\ker(u)\oplus\text{Im}(u)$ ;
		\item il existe une base de $E$ dans laquelle la matrice de $u$ est de la forme $\displaystyle\begin{pmatrix}
			A&0\\
			0&0
		\end{pmatrix}$ avec $A$ inversible.
	\end{enumerate}
\end{ex}

\begin{ex}[Moulin.Mat.34]{$\star$ Similitude et surcorps}
	Montrer que si $A$ est $B$ dans $\mathcal{M}_n(\R)$ sont semblables dans $\mathcal{M}_n(\C)$, alors elles sont aussi semblables dans $\mathcal{M}_n(\R)$. Même question (mais c'est plus dur) avec $\K$ et $\L$ un surcorps de $\K$ (on pourra commencer par le cas où $\L$ est de dimension finie sur $\K$).
\end{ex}
\begin{proof}
	On commence par le cas de $\R$ et $\C$ qui est le cas "usuel".
	\begin{itemize}
		\item On fixe $P\in\mathcal{GL}_n(\C)$ telle que $PB=AP$. On écrit $P=P_1+iP_2$ avec $P_1$ et $P_2$ à coefficients réels. On obtient alors : $$\begin{cases}
			P_1B=AP_1\\
			P_2B=AP_2
		\end{cases}$$ Pour $\lambda\in\C$, on pose $P_\lambda=P_1+\lambda P_2$ et on veut montrer qu'il existe un $\lambda$ réel tel que $P_\lambda$ soit inversible. Par l'absurde, si ce n'était pas le cas, le déterminant de $P_\lambda$, polynomial en $\lambda$, serait nul sur $\R$ qui est infini donc sur $\C$ ce qui est absurde car il n'est pas nul en $i$.
		\item On suppose $\dim_\K(\L)=n<+\infty$ et on écrit $P=l_1P_1+\ldots+l_nP_n$ avant de considérer $$P=\lambda_1P_1+\ldots+\lambda_nP_n$$ et $\varphi$ qui à $(\lambda_1,\ldots,\lambda_n)$ associe le déterminant de $P$. C'est un polynôme à $n$ indéterminées. On montre alors par récurrence sur $n$ que $$\forall Q\in\L[X_1,\ldots,X_n],\ (\forall(x_1,\ldots,x_n)\in\K^n,\ Q(x_1,\ldots,x_n)=0)\implies Q=0$$ OK pour l'initialisation, pour l'hérédité, on fixe $x_0\in\K$ et par hypothèse de récurrence $$Q(x_0,X_1,\ldots,X_n)=0$$ Puis on a donc $$Q(X_0,x_1,\ldots,x_n)$$ qui s'annule en une infinité de points (tous les $x_0$) donc $Q(X_0,x_1,\ldots,x_n)=0$ quel que soit le $n-uplet$. On en déduit que $Q(X_0,\ldots,X_n)$ est le polynôme nul. Avec ce lemme, on conclut de même que dans le premier point.
		\item Dans le cas général, on considère $F=\text{Vect}(P_{i,j})$ puis on utilise le cas précédent avec une base de $F$.
	\end{itemize}
\end{proof}

\begin{ex}[Moulin.Mat.35]{$\circ$}
	Soit $u$ un endomorphisme d'un espace vectoriel $E$ de dimension $3$. montrer que si $u\ne 0$ et $u^2=0$, il existe une base de $E$ dans laquelle la matrice de $u$ est $$\begin{pmatrix}
		0&1&0\\
		0&0&0\\
		0&0&0
	\end{pmatrix}$$
\end{ex}

\begin{ex}[Moulin.Mat.36]{}
	Soit $E$ un $\K$-espace vectoriel de dimension finie $n\geqslant 1$. Soit $u\in\mathcal{L}(E)$ non nul tel que $u^2=0$.
	\begin{enumerate}
		\item Montrer l'existence d'une base de $E$ dans laquelle la matrice de $u$ est diagonale par blocs $\text{Diag}(\underbrace{A,\ldots,A}_{r},0_{n-2r})$ où $r$ est le rang de $u$ et $$A=\begin{pmatrix}
			0&1\\
			0&0
		\end{pmatrix}\in\mathcal{M}_2(\K)$$
		\item Montrer que $u$ est la différence de deux matrices de projecteurs.
		\item La similitude induit une relation d'équivalence sur l'ensemble des matrices $M$ de $\mathcal{M}_n(\K)$ telles que $M^2=0$. Combien y a-t-il de classes d'équivalence ?
	\end{enumerate}
\end{ex}

\begin{ex}[Moulin.Mat.37]{}
	Soit $A\in\mathcal{M}_n(\K)$, $B\in\mathcal{M}_p(\K)$ et $C\in\mathcal{M}_{n,p}(\K)$. On considère les assertions suivantes :
	\begin{enumerate}[label=(\roman*)]
		\item $\exists X\in\mathcal{M}_{n,p}(\K),\ AX-XB=C$ ;
		\item $\displaystyle\begin{pmatrix}
			A&0\\
			0&B
		\end{pmatrix}$ et $\displaystyle\begin{pmatrix}
		A&C\\
		0&B
		\end{pmatrix}$ sont semblables.
	\end{enumerate}
	\begin{enumerate}
		\item Montrer que $(i)$ implique $(ii)$.
		\item $\star\star$ En étudiant les noyaux de : $$\varphi_0:M\in\mathcal{M}_{n+p}(\K)\mapsto \begin{pmatrix}
			A&0\\
			0&B
		\end{pmatrix}M-M\begin{pmatrix}
		A&0\\
		0&B
		\end{pmatrix},$$
		$$\varphi_0:M\in\mathcal{M}_{n+p}(\K)\mapsto \begin{pmatrix}
			A&C\\
			0&B
		\end{pmatrix}M-M\begin{pmatrix}
			A&0\\
			0&B
		\end{pmatrix},$$ montrer que $(ii)$ implique $(i)$.
	\end{enumerate}
\end{ex}

\section{Trace}

\begin{ex}[Moulin.Mat.38]{$\circ$}
	Soit $A\in\mathcal{M}_{n,p}(\K)$ et $B\in\mathcal{M}_{p,n}(\K)$. Montrer que $(AB)^k$ et $(BA)^k$ ont même trace ($k\in$ ?).
\end{ex}

\begin{ex}[Moulin.Mat.39]{$\Delta$ Représentation des formes linéaires sur $\mathcal{M}_n(\K)$}
	Soit $\psi$ une forme linéaire sur $\mathcal{M}_n(\K)$. Montrer qu'il existe une unique matrice $A\in\mathcal{M}_n(\K)$ telle que $$\forall M\in\mathcal{M}_n(\K),\ \psi(M)=\text{Tr(AM)}$$
\end{ex}
\begin{proof}
	Considérer pour tout $A\in\mathcal{M}_n(\K)$ l'application $$\begin{array}{lrcl}
		f_A:&\mathcal{M}_n(\K)&\longrightarrow&\K \\
		&M&\longmapsto&\Tr(AM)
	\end{array}$$ puis considérer $$\begin{array}{lrcl}
		\varphi:&\mathcal{M}_n(\K)&\longrightarrow&\mathcal{M}_n(\K)\st\\
		&A&\longmapsto&f_A
	\end{array}$$ Déjà, il est clair que les $f_A$ sont bien des formes linéaires et que $\varphi$ est bien définie et linéaire. De plus, on montre que $\varphi$ est injective car si $\forall M\in\mathcal{M}_n(\K),\ \Tr(AM)=0$, alors en utilisant des matrices d'opérations élémentaires et en annulant certaines lignes ou colonnes, on montre que $A$ est nulle. Si $\K=\R$, il est plus rapide de le faire avec $M=A^T$, et si $\K=\C$, il est plus rapide de le faire avec $M=\overline{A}^T$ (cela correspond aux produits scalaires). Par égalité de dimension, $\varphi$ est donc un isomorphisme, ce qui conclut.
\end{proof}

\begin{ex}[Moulin.Mat.40]{}
	En utilisant l'exercice précédent :
	\begin{enumerate}
		\item Montrer que tout hyperplan de $\mathcal{M}_n(\K)$ contient une matrice inversible pour $n>1$.
		\item Trouver toutes les formes linéaires $f$ sur $\mathcal{M}_n(\K)$ vérifiant : $$\forall (X,Y)\in\mathcal{M}_n(\K)^2,\ f(XY)=f(YX)$$
	\end{enumerate}
\end{ex}
\begin{proof}
	Applications classiques de l'exercice précédent.
	\begin{enumerate}
		\item On prend la forme linéaire non nulle dont $H$ est le noyau. On représente cette forme linéaire par $M\mapsto\Tr(AM)$ avec $A$ non nulle. Par équivalence, on se ramène à chercher une matrice inversible $N$ telle que $\Tr(J_r N)=0$. On pourra considérer pour $N$ la matrice du cycle $(1\ 2\ \ldots \ n)$.
		\item 
	\end{enumerate}
\end{proof}

\begin{ex}[Moulin.Mat.41]{$\star\star$}
	Existe-t-il dans $\mathcal{M}_n(\K)$ un hyperplan $H$ stable par multiplication ? Indication : lorsque $n\geqslant3$, on pourra utiliser la représentation des formes linéaires pour trouver une matrice $A$ telle que $\forall X\in H,\ AX\in\text{Vect}(A)$.
\end{ex}

\begin{ex}[Moulin.Mat.42]{$\star\star$}
	Déterminer les hyperplans $H$ de $\mathcal{M}_n(\C)$ stables par conjugaison, c'est-à-dire tels que : $$\forall P\in\mathcal{GL}_n(\C),\ \forall M\in H,\ PMP\inv\in H$$
\end{ex}

\begin{ex}[Moulin.Mat.43]{}
	Déterminer les formes linéaires $f$ sur $\mathcal{M}_n(\K)$ vérifiant $f(ABC)=f(CBA)$ quelles que soient les matrices $A$, $B$ et $C$ de $\mathcal{M}_n(\K)$.
\end{ex}

\begin{ex}[Moulin.Mat.44]{$\star$ Matrices de trace nulle}
	\begin{enumerate}
		\item Montrer qu'une matrice de $\mathcal{M}_n(\K)$ est de trace nulle si, et seulement si, elle est semblable à une matrice dont tous les termes diagonaux sont nuls. On pourra procéder par récurrence.
		\item Montrer que toute matrice de trace nulle est somme de deux matrices nilpotentes.
		\item Montrer qu'une matrice de $\mathcal{M}_n(\K)$ est de trace nulle si, et seulement si, elle s'écrit $XY-YX$ avec $(X,Y)\in\mathcal{M}_n(\K)^2$.
	\end{enumerate}
\end{ex}

\begin{ex}[Moulin.Mat.45]{$\star\star$}
	Montrer que pour tout $(a_1,\ldots,a_n)\in\K^n$ tel que $\Tr(A)=a_1+\cdots+a_n$, il existe une matrice semblable à $A$ dont les coefficients diagonaux sont $a_1,\ldots,a_n$.
\end{ex}

\begin{ex}[Moulin.Mat.46]{Sommes de symétries dans $\mathcal{M}_n(\C)$}
	\begin{enumerate}
		\item Que peut-on dire de la trace d'une somme de matrices de symétries ?
		\item Montrer que $\displaystyle\begin{pmatrix}
			-1&a\\
			0&1
		\end{pmatrix}$ est une matrice de symétrie.
		\item Montrer que toute matrice de trace nulle est semblable à une matrice $(m_{i,j})$ avec $m_{1,1}=0$.
		\item $\star\star$ En déduire la réciproque de la première question lorsque $n=2$, puis pour $n$ quelconque.
	\end{enumerate}
\end{ex}

\begin{ex}[Moulin.Mat.47]{}
	\begin{enumerate}
		\item Soit $M\in\mathcal{M}_n(\C)$. Résoudre $X+X^T=\Tr(X)M$ dans $\mathcal{M}_n(\C)$.
		\item Soit $(A,B)\in\mathcal{M}_n(\K)^2$. Résoudre $X+\Tr(X)A=B$.
	\end{enumerate}
\end{ex}

\begin{ex}[Moulin.Mat.48]{}
	Montrer que tout automorphisme de l'algèbre $\mathcal{M}_n(\K)$ conserve la trace.
\end{ex}

\section{Systèmes linéaires, opérations élémentaires}

\begin{ex}[Moulin.Mat.49]{$\Delta$ Matrices à diagonale dominante}
	On dit qu'une matrice $A=(a_{i,j})\in\mathcal{M}_n(\C)$ est à diagonale dominante lorsque : $$\forall i\in\llbracket 1,n\rrbracket,\ \lvert a_{i,i}\rvert>\sum_{j\ne i}\lvert a_{i,j}\rvert$$
	\begin{enumerate}
		\item $\Delta$ Montrer qu'une matrice à diagonale dominante est inversible. On pourra raisonner par l'absurde et considérer $X\ne 0$ tel que $AX=0$. Il peut être instructif de redémontrer ce résultat par récurrence sur $n$ à l'aide d'un pivot de Gauss.
		\item En déduire une borne sur l'ensemble des valeurs propres d'une matrice $A\in\mathcal{M}_n(\C)$ quelconque.
	\end{enumerate}
\end{ex}
\begin{proof}
	\begin{enumerate}
		\item On raisonne par l'absurde en considérant un élément non nul $X$ du noyau de $A$. On écrit la relation $AX=0$ puis on s'intéresse à la ligne de l'égalité précédente qui correspond au $x_i$ de module maximal non nul et on obtient une absurdité grâce au fait que la diagonale est dominante.
		\item Disque de Gershgoring.
	\end{enumerate}
\end{proof}

\begin{ex}[Moulin.Mat.50]{$\Delta$}
	Résoudre le système linéaire : $$\forall i\in\llbracket0,n\rrbracket,\ \sum_{j=0}^{n}a_i^jx_j=a_i^{n+1}$$ où $(a_0,\ldots,a_n)$ est une famille fixée de scalaires distincts deux à deux.
\end{ex}

\begin{ex}[Moulin.Mat.51]{$\Delta$}
	Soit $$A=\begin{pmatrix}
		1+\lambda^2&-\lambda&&(0)\\
		-\lambda&1+\lambda^2&\ddots&\\
		&\ddots&\ddots&-\lambda\\
		(0)&&-\lambda&1+\lambda^2
	\end{pmatrix}$$ avec $\lambda\in\C$. Résoudre le système $AX=0$. On pourra introduire une récurrence linéaire.
\end{ex}

\begin{ex}[Moulin.Mat.52]{}
	Résoudre : $$\forall i\in\llbracket0,n\rrbracket,\ \sum_{j=0}^{n-1}j^ix_j=\lambda^i$$
\end{ex}

\begin{ex}[]{$\Delta$ Générateurs de $\mathcal{SL}_n(\K)$ et de $\mathcal{GL}_n(\K)$}
	Montrer que les matrices de transvection engendrent $\mathcal{SL}_n(\K)$. Montrer que toute matrice de $\mathcal{GL}_n(\K)$ s'écrit comme le produit de matrices de transvection et d'une matrice de dilatation.
\end{ex}
\begin{proof}
	On réalise un pivot de Gauss en utilisant uniquement des transvections.
	\begin{itemize}
		\item On se ramène à $a_{1,1}=1$ par la transvection $L_1\leftarrow\displaystyle L_1+\frac{1-a_{1,1}}{a_{k,1}}L_k$ avec $k$ tel que $a_{k,1}\ne 0$.
		\item Ensuite, en effectuant les transvections $L_k\leftarrow L_k-a_{k,1}L_1$ pour $k\in\llbracket2,n\rrbracket$, on se ramène à : $$T_{d_1}\cdots T_1A=\begin{pmatrix}
			1&*&*&*\\
			0&&&\\
			\vdots&&(*)&\\
			0&&&
		\end{pmatrix}$$
		\item On se ramène à $a_{1,1}=1$ par la transvection $L_1\leftarrow\displaystyle L_1+\frac{1-a_{1,1}}{a_{k,1}}L_k$ avec $k$ tel que $a_{k,1}\ne 0$.
		\item Puis, en effectuant les transvections $C_k\leftarrow C_k-a_{1,k}C_1$ pour $k\in\llbracket2,n\rrbracket$, on se ramène à : $$T_{d_1}\cdots T_1A\tilde{T}_{1}\cdots\tilde{T}_{e_1}=\begin{pmatrix}
			1&0&\cdots&0\\
			0&&&\\
			\vdots&&A_1&\\
			0&&&
		\end{pmatrix}$$ avec $A_1\in\mathcal{GL}_{n-1}(\K)$.
		\item On réitère sur $A_1$ par récurrence pour obtenir :  $$T_{d_1}\cdots T_1A\tilde{T}_{1}\cdots\tilde{T}_{e_1}=\text{Diag}(1,\cdots,1,k)$$ avec $k=\det(A)$.
	\end{itemize}
	Ensuite, si $A\in\mathcal{SL}_n(\K)$, $k=\det(A)=1$ donc on obtient le résultat. Et sinon, $k\ne 0$ donc on obtient le résultat pour $\mathcal{GL}_n(\K)$.
\end{proof}

\section{Déterminants}

\begin{ex}[Moulin.Mat.53]{$\Delta$}
	Soit $E$ un espace vectoriel de dimension $n$ et $u\in\mathcal{L}(E)$. Montrer que pour toute forme $n$-linéaire alternée $f$ sur $E$, on a : $$\forall (x_1,\ldots,x_n)\in E^n,\ \sum_{i=1}^{n}f(x_1,\ldots,x_{i-1},u(x_i),x_{i+1},\ldots,x_n)=\Tr(u)f(x_1,\ldots,x_n)$$
\end{ex}

\begin{ex}[Moulin.Mat.54]{$\Delta$}
	Montrer que l'espace vectoriel des formes $p$-linéaires alternées sur un espace vectoriel de dimension $n$ est de dimension $\displaystyle\binom{n}{p}$.
\end{ex}

\begin{ex}[Moulin.Mat.55]{}
	Calculer le déterminant $\displaystyle\begin{vmatrix}
		1&2&\cdots&n\\
		n&1&\cdots&n-1\\
		&\vdots&&\vdots\\
		2&\cdots&n&1
	\end{vmatrix}$.
\end{ex}

\begin{ex}[Moulin.Mat.56]{$\Delta$ Vers la densité des matrices inversibles}
	Soit $A\in\mathcal{M}_n(\K)$ où $\K$ est un corps infini. Montrer qu'il existe une infinité de $\lambda\in\K$ tels que $\lambda I_n-A$ soit inversible.
\end{ex}
\begin{proof}
	L'application qui à $\lambda$ associe le déterminant de $\lambda I_n-A$ est polynomiale en $\lambda$. Elle est non nulle car de degré $n\geqslant 1$ (sinon l'exercice n'a pas d'intérêt). donc elle possède un nombre fini de racines, c'est-à-dire qu'il existe une infinité de $\lambda$ tels que le déterminant de $\lambda I_n-1$ soit non nul, \textit{ie} que $\lambda I_n-A$ soit inversible.
\end{proof}

\begin{ex}[Moulin.Mat.57]{$\circ$}
	Soit $\sigma\in\mathcal{S}_n$. Calculer le déterminant de la matrice $P_\sigma=\left(\delta_{i,\sigma(j)}\right)_{1\leqslant i,j\leqslant n}$.
\end{ex}

\begin{ex}[Moulin.Mat.58]{}
	Soit $p\in\llbracket0,n-1\rrbracket$. Calculer le déterminant : $$\begin{vmatrix}
		a_1^0&\cdots&a_1^{p-1}&a_1^{p+1}&\cdots&a_1^n\\
		\vdots&&\vdots&\vdots&&\vdots\\
		a_p^0&\cdots&a_p^{p-1}&a_p^{p+1}&\cdots&a_p^n\\
		\vdots&&\vdots&\vdots&&\vdots\\
		a_n^0&\cdots&a_n^{p-1}&a_n^{p+1}&\cdots&a_n^n\\
	\end{vmatrix}$$
\end{ex}

\begin{ex}[Moulin.Mat.59]{Déterminant circulant}
	Soit $(a_0,\ldots,a_{n-1})\in\C^n$. Calculer le déterminant de la matrice $(b_{i,j})_{1\leqslant i,j\leqslant n}$ avec $$b_{i,j}=a_{[j-i]}$$ où $[k]$ est le reste de la division euclidienne de $k$ par $n$. On pourra utiliser une matrice de Vandermonde associée aux racines $n$-èmes de l'unité. \\
	Application : factoriser $a^3+b^3+c^3-3abc$ dans $\C$, puis dans $\R$.
\end{ex}

\begin{ex}[Moulin.Mat.60]{Déterminant de Cauchy}
	Calculer le déterminant de Cauchy : $$C(a_1,\ldots,a_n,b_1,\ldots,b_n)=\displaystyle\begin{vmatrix}
		\displaystyle\frac{1}{a_1+b_1}&\displaystyle\frac{1}{a_1+b_2}&\cdots&\displaystyle\frac{1}{a_1+b_n}\\
		\displaystyle\frac{1}{a_2+b_1}&\displaystyle\frac{1}{a_2+b_2}&\cdots&\displaystyle\frac{1}{a_2+b_n}\\
		\vdots&\vdots&&\vdots\\
		\displaystyle\frac{1}{a_n+b_1}&\displaystyle\frac{1}{a_n+b_2}&\cdots&\displaystyle\frac{1}{a_n+b_n}
	\end{vmatrix}$$ où $(a_1,\ldots,a_n)$ et $(b_1,\ldots,b_n)$ sont des familles de complexes tels que $$\forall (i,j)\in\llbracket1,n\rrbracket^2,\ a_i+b_\ne 0$$
\end{ex}
\begin{proof}
	Montrer par récurrence la formule $$C_n=\frac{V(a_1,\ldots,a_n)V(b_1,\ldots,b_n)}{\prod_{i,j}(a_i+b_j)}$$On multiplie toutes les colonnes par $a_n+b_j$, puis on soustrait la dernière colonne aux autres avant de factoriser les termes communs aux lignes et aux colonnes pour finalement développer par rapport à la dernière ligne. On obtient la relation de récurrence : $$C_n=\frac{1}{a_n+b_n}\left(\prod_{i=1}^{n-1}\frac{(a_n-a_i)(b_n-b_i)}{(a_n+b_i)(a_i+b_n)}\right)C_{n-1}$$
\end{proof}
\begin{app}
	Calcul du déterminant de la matrice de Hilbert.
\end{app}

\begin{ex}[Moulin.Mat.61]{}
	Soit $a_1,\ldots,a_n,b_1,\ldots,b_n$ des nombres complexes. Calculer le déterminant : $$\begin{vmatrix}
		a_1+b_1&b_1&\cdots&b_1\\
		b_2&a_2+b_2&\cdots&b_2\\
		\vdots&&\ddots&\\
		b_n&\cdots&b_n&a_n+b_n
	\end{vmatrix}$$
\end{ex}
\begin{proof}
	Utiliser la multilinéarité.
\end{proof}

\begin{ex}[Moulin.Mat.62]{}
	Soit $a_1,\ldots,a_n,b_1,\ldots,b_n$ des nombres complexes.
	\begin{enumerate}
		\item Calculer le déterminant de la matrice $((a_i+b_j)^{n-1})_{1\leqslant i,j\leqslant n}$. On pourra penser à un produit matriciel.
		\item Calculer le déterminant de la matrice $((a_i+b_j)^{p})_{1\leqslant i,j\leqslant n}$ pour $p\in\llbracket 0,n-2\rrbracket$.
	\end{enumerate}
\end{ex}

\begin{ex}[Moulin.Mat.63]{$\Delta$}
	Soit $n \in \N\st$.
	Soit $M \in \mathcal{M}_n(\Z)$. Montrer que $$ M \in \mathcal{GL}_n(\Z) \Longleftrightarrow \det(M)=\pm 1 $$ où on a posé $$\mathcal{GL}_n(\Z)=\left\{A \in \mathcal{M}_n(\Z) \mid \exists B \in \mathcal{M}_n(\Z), \ AB=I_n\right\}$$
\end{ex}
\begin{proof}
	Pour $\Longrightarrow$, utiliser le fait que le déterminant de $M$ doit être une unité de $\Z$. Pour $\Longleftarrow$, utiliser la formule de la comatrice et le fait que le déterminant est polynomial.
\end{proof}

\begin{ex}[Moulin.Mat.64]{}
	Calculer le déterminant de la matrice de $\mathcal{M}_n(\R)$ de terme général $\sin(a_i+b_j)$.
\end{ex}

\begin{ex}[Moulin.Mat.65]{}
	Calculer le déterminant de la matrice de $\mathcal{M}_n(\R)$ de terme général $\lvert j-i\rvert$.
\end{ex}

\begin{ex}[Moulin.Mat.66]{}
	Calculer le déterminant $\displaystyle\begin{vmatrix}
		1+x^2&x&&(0)\\
		x&1+x^2&\ddots&\\
		&\ddots&\ddots&x\\
		(0)&&x&1+x^2
	\end{vmatrix}$.
\end{ex}

\begin{ex}[Moulin.Mat.67]{}
	Calculer le déterminant $\displaystyle\begin{vmatrix}
		a+b&ab&&(0)\\
		1&a+b&\ddots&\\
		&\ddots&\ddots&ab\\
		(0)&&1&a+b
	\end{vmatrix}$.
\end{ex}

\begin{ex}[Moulin.Mat.68]{$\star$ Produit de Kronecker} Pour $A=(a_{i,j})\in\mathcal{M}_n(\K)$ et $B\in\mathcal{M}_m(\K)$, on note : $$A\otimes B=\begin{pmatrix}
		a_{1,1}B&a_{1,2}B&\cdots&a_{1,n}B\\
		a_{2,1}B&a_{2,2}B&\cdots&a_{2,n}B\\
		\vdots&\vdots&&\vdots\\
		a_{n,1}B&a_{n,2}B&\cdots&a_{n,n}B
	\end{pmatrix}\in\mathcal{M}_{mn}(\K)$$
	\begin{enumerate}
		\item Montrer que $\det(A\otimes B)=\det(A)^m\det(B)^n$.
		\item En déduire le déterminant de l'application de $\mathcal{M}_{n,m}(\K)$ dans $\mathcal{M}_{n,m}(\K)$ qui à $M$ associe $AMB$.
	\end{enumerate}
\end{ex}\begin{proof}
	En effet, on utilise la relation fonctionnelle $(AC)\otimes(BD)=(A\otimes B)\times(C\otimes D)$ avec des matrices identités. Le plus dur est alors de montrer que $A\otimes B$ est semblable à $B\otimes A$, ce qui permettra de conclure. Pour cela, prendre une base adaptée à la structure de produit tensoriel, et, au lieu de sommer d'abord sur les indices des premières composantes de cette famille produit tensoriel, sommer sur l'autre composante. Cela permet de passer de $A\otimes B$ à $B\otimes A$ en changeant de base, donc ces matrices sont semblables et ont le même déterminant. On conclut en appliquant ceci avec les matrices identités car le déterminant de $I\otimes B$ se calcule très bien comme déterminant diagonal par blocs.
\end{proof}

\begin{ex}[Moulin.Mat.69]{}
	Calculer le déterminant de l'endomorphisme qui à une matrice associe sa transposée.
\end{ex}
\begin{proof}
	L'essentiel est de mettre la base canonique de $\mathcal{M}_n(\K)$ dans le bon ordre. On met d'abord les $E_{i,i}$, puis on met ensemble les $E_{i,j}$ et $E_{j,i}$. On trouve en calculant le déterminant par blocs dans cette base : $$(-1)^{n(n-1)/2}$$
\end{proof}

\begin{ex}[Moulin.Mat.70]{$\circ$}
	Soit $A\in\mathcal{M}_{p,n}(\K)$ et $B\in\mathcal{M}_{n,p}(\K)$. En effectuant dans les deux sens le produit par blocs entre les matrices : $$\begin{pmatrix}
		I_n&B\\
		A&\lambda I_p
	\end{pmatrix}\ \ \text{et} \ \ \begin{pmatrix}
		\lambda I_n&0\\
		-A&I_p
	\end{pmatrix}$$ montrer que $$\forall\lambda\in\K,\ \lambda^n\det(\lambda I_p-AB)=\lambda^p\det(\lambda I_n-BA)$$
\end{ex}
\begin{proof}
	Effectuer les deux produits par blocs. On trouve des matrices triangulaires par blocs. Calculer les deux déterminants, qui sont égaux par commutativité de $\K$. On trouve bien la formule demandée.
\end{proof}
\begin{app}
	Permet de montrer que $\chi_{AB}=\chi_{BA}$ en prenant $p=n$.
\end{app}

\begin{ex}[Moulin.Mat.71]{}
	Soit $(A,B,C,D)\in\mathcal{M}_n(\K)^4$ tel que $CD=CD$. Montre que $$\det\begin{pmatrix}
		A&B\\
		C&D
	\end{pmatrix}=\det(AD-BC)$$ On pourra commencer par le cas où $D$ est inversible.
\end{ex}
\begin{proof}
	Effecteur des transvections par blocs en multipliant :
	\begin{itemize}
		\item à droite par $\displaystyle\begin{pmatrix}
			D&0\\
			-C&I_n
		\end{pmatrix}$
		\item à gauche par $\displaystyle\begin{pmatrix}
			I_n&0\\
			0&D\inv
		\end{pmatrix}$
	\end{itemize} puis calculer le déterminant obtenu. Dans le cas général, utiliser la densité des matrices inversibles (en considérant $D_X=D-X.I_n$ par exemple).
\end{proof}

\begin{ex}[Moulin.Mat.72]{}
	Soit $z_1,\ldots,z_n$ les racines complexes du polynôme : $$X^n+a_{n-1}X^{n-1}+\cdots+a_1X+a_0$$ Montrer : $$\begin{vmatrix}
		n&z_1&\cdot&z_n \\
		z_1&1+z_1^2&&\\
		\vdots&(1)&\ddots&(1)\\
		z_n&&&1+z_n^2
	\end{vmatrix}=a_1^2$$ Indication : on pourra utiliser la somme des inverses des $z_i$ lorsqu'ils sont tous non nuls.
\end{ex}

\begin{ex}[Moulin.Mat.73]{$\Delta\star$}
	Soit $A\in\mathcal{M}_n(\Z)$ ayant tous les coefficients diagonaux impairs et tous les coefficients non diagonaux pairs. Montrer que $A$ est inversible dans $\mathcal{M}_n(\R)$.
\end{ex}
\begin{proof}
	Déjà, $\det(A)\in\Z$ et, si on note $\overline{A}$ la matrice des classes des coefficients de $A$ dans $\Z/2\Z$,  la formule des signatures prouve que $$\overline{\det(A)}=\det(\overline{A})$$ Or, $\overline{A}=I_n$ donc $\det(\overline{A})=\overline{1}$. Par conséquent, $\det(A)$ est impair donc non nul donc $A$ est inversible dans $\mathcal{M}_n(\R)$.
\end{proof}

\begin{ex}[Moulin.Mat.74]{$\star$}
	Soit $A$ et $B$ dans $\mathcal{M}_n(\R)$ telles que $A^2+B^2=\sqrt{3}(AB-BA)$ et $AB-BA$ soit inversible. Montrer que $n$ est un multiple de $6$.
\end{ex}
\begin{proof}
	L'astuce est de calculer $(A+iB)(A-iB)$. L'hypothèse nous donne : $$(A+iB)(A-iB)=(\sqrt{3}-i)(AB-BA)=2e^{-i\pi/6}(AB-BA)$$ Or, la formule des signatures donne rapidement $$\overline{\det(A+iB)}=\det(A-iB)$$ Par conséquent, $\det((A+iB)(A-iB))\in\R$. Or, $$\det((A+iB)(A-iB))=2^ne^{-in\pi/6}\underbrace{\det(AB-BA)}_{\in\R\st}$$ Par conséquent, $e^{-in\pi/6}\in\R$ donc $6$ divise $n$.
\end{proof}

\begin{ex}[Moulin.Mat.75]{}
	Soit $M$ une matrice inversible de $\mathcal{M}_n(\C)$ que l'on décompose par blocs en : $$M=\begin{pmatrix}
		A&B\\
		C&D
	\end{pmatrix}$$ avec $A$ et $D$ carrées. On décompose, avec le même format, la matrice $$M\inv=\begin{pmatrix}
	A'&B'\\
	C'&D'
	\end{pmatrix}$$ Montrer qu'alors $\det(A)=\det(M)\det(D')$. Indication : dans le cas où $A$ est inversible, on pourra commencer par exprimer $D'$ à l'aide de $A$, $B$, $C$ et $D$.
\end{ex}

\begin{ex}[Moulin.Mat.76]{}
	Soit $(A,B)\in\mathcal{M}_n(\K)^2$ et la matrice $M\in\mathcal{M}_{2n}(\K)$ définie par blocs par $$M=\begin{pmatrix}
		A&B\\
		B&A
	\end{pmatrix}$$ Montrer alors que $\det(M)=\det(A+B)\det(A-B)$.
\end{ex}
\begin{proof}
	Effectuer les transvections par blocs suivantes :
	\begin{itemize}
		\item $C_1\leftarrow C_1-C_2$, codée par une multiplication à droite par $$\begin{pmatrix}
			I_n&0\\
			-I_n&I_n
		\end{pmatrix}$$
		\item puis $L_2\leftarrow L_1+L_2$ codée par une multiplication à gauche par $$\begin{pmatrix}
			I_n&0\\
			I_n&I_n
		\end{pmatrix}$$
	\end{itemize} Calculer ensuite le déterminant des matrices obtenues.
\end{proof}

\begin{ex}{Déterminant et matrice de rang 1}
	Soit $n \in \N\st$ et $A \in \mathcal{M}_n(\K)$ de rang 1. Montrer que $$\text{det}(I_n+A) = 1 + \Tr(A)$$
\end{ex}
\begin{proof}
	Comme la matrice est de rang 1, elle possède une colonne non nulle et on peut exprimer toutes les autres colonnes en fonction de cette colonne principale. Alors, on peut développer le déterminant souhaité par multilinéarité, et à ce moment là, beaucoup de termes disparaissent par le caractère alterné du déterminant. Après quelques calculs, on obtient le résultat voulu.
\end{proof}

\begin{ex}{Déterminants en vrac}
	Soit $n \in \N\st$. Calculer
	\begin{enumerate}
		\item $\det((i^j)_{1\le i,j \le n})$
		\item $\det((\max(i,j))_{1\le i,j \le n}$)
	\end{enumerate}
\end{ex}
\begin{proof}
	Représenter le début de la matrice peut aider...
	\begin{enumerate}
		\item Faire sortir un $i$ de chaque ligne par multilinéarité. On obtient alors $n!$ fois un déterminant de Vandermonde qui vaut $\prod_{1\le i < j \le n} (j-i)=\prod_{j=2}^{n} \prod_{i=1}^{j-1} (j-i)=\prod_{j=1}^{n-1} j!$. On en tire le résultat final : $$ \det((i^j)_{1\le i,j \le n})= \prod_{k=1}^{n} k!$$
		\item C'est surtout ici que la représentation de la matrice aide : la matrice est "en escalier". Cela peut aider à savoir quelles transvections faire. On pourra formaliser avec une récurrence si on le souhaite.
	\end{enumerate}
\end{proof}
\begin{ex}{$\Delta\star$ Déterminant de Smith} Calculer le déterminant de Smith :$$\det\left(i\wedge j\right)_{(1\leqslant i,j\leqslant n)}$$
\end{ex} \begin{proof}
	Utiliser la formule $$n=\sum_{d\mid n}\varphi(d)$$ Introduire la notation $\delta_{i\mid j}$ qui vaut $1$ si $i$ divise $j$ et $0$ sinon. Ainsi, on peut interpréter $i\wedge j$ comme le terme d'un produit matriciel. Dessiner les matrices correspondantes, dont on montre qu'elles sont triangulaires à cause des $\delta_{i\mid j}$. On obtient : $$\det\left(i\wedge j\right)_{(1\leqslant i,j\leqslant n)}=\prod_{k=1}^{n}\varphi(k)$$
\end{proof}

\begin{ex}{$\Delta\star$ Le problème des $2+1$ vaches/cailloux}
	Voici la généralisation du célèbre problème des 11 vaches. Soit $n \ge 2$.
	\begin{enumerate}
		\item Soit $M \in \mathcal{M}_n(\Z)$ dont tous les coefficients valent $\pm1$ sauf ceux sur la diagonale qui valent 0. Montrer que $\det(M) \ne 0$.
		\item Soit un groupe de $2n+1$ vaches (ou cailloux, au choix) tel que dès que l'on retire une vache, on peut toujours former parmi les $2n$ restantes deux groupes de $n$ vaches dont les sommes des poids des vaches par groupes sont égales. Montrer que toutes les vaches ont le même poids. Il est conseillé de traduire matriciellement le problème, la première question n'est pas là pour rien... On pourra montrer que l'image de la matrice considérée est de dimension $2n$.
	\end{enumerate}
\end{ex}
\begin{proof}
	Si on n'a pas de chance, c'est sans la première question !
	\begin{enumerate}
		\item Passer dans $\Z/2\Z$
		\item Montrer que la matrice considérée est de rang supérieur à $2n$ en utilisant la question précédente, puis montrer que le noyau est de dimension supérieure ou égale à 1 car il contient le vecteur colonne composé uniquement de 1. Conclure en remarquant que le vecteur "poids des vaches" est dans le noyau. Il s'écrit donc comme une dilatation du vecteur colonne comportant uniquement des 1 : toutes les vaches ont donc le même poids.
	\end{enumerate}
\end{proof}
\begin{ex}{$\star\star$ Une somme étrange (pas classique du tout mais marrant)}
	Soit $n\geqslant1$. Pour tout $\sigma\in\mathfrak{S}_n$, on note $\text{inv}(\sigma)$ le nombre de points de $\llbracket1,n\rrbracket$ invariants par $\sigma$. Calculer : $$\sum_{\sigma\in\mathfrak{S}_n}\frac{\varepsilon(\sigma)}{\text{inv}(\sigma)+1}$$
\end{ex}
\begin{proof}
	Écrire $$\displaystyle\frac{1}{\text{inv}(\sigma)+1}=\int_{0}^{1}x^{\text{inv}(\sigma)}\ dx$$ puis utiliser la linéarité de l'intégrale. S'intéresser au polynôme obtenu sous le symbole d'intégration. Interpréter ce polynôme comme un déterminant (suite de l'indication : celui de la matrice comportant des $1$ partout et une diagonale de $x$). Calculer ensuite ce déterminant par opérations.
\end{proof}

\begin{ex}[Châteaux]{Déterminant de Pascal}
	Pour tout $n\in\N$, on note $B_n\in\mathcal{M}_{n+1}(\R)$ de terme général : $$[B_n]_{i,j}=\binom{i+j}{j}$$ où $(i,j)\in\llbracket0,n\rrbracket$ (on numérote les lignes et les colonnes à partir du rang $0$ pour avoir des formules plus jolies). On se propose de démontrer $\forall n\in\N,\ \det(B_n)=1$ de deux façons différentes.
	\begin{enumerate}
		\item \begin{enumerate}[label=\alph*.]
			\item Soit $(x_0,\ldots,x_n)\in\R^{n+1}$. Soit $(P_0,\ldots,P_n)\in\R[X]^{n+1}$ tels que, pour tout $k\in\llbracket0,n\rrbracket$, le polynôme $P_k$ soit de degré $k$ et de coefficient dominant $\lambda_k$. Calculer le déterminant de la matrice $\left(P_i(x_j)\right)_{(i,j)\in\llbracket0,n\rrbracket^2}$.
			\item Conclure. 
		\end{enumerate}
		\item \begin{enumerate}[label=\alph*.]
			\item Montrer : $\displaystyle\forall(i,j)\in\N^2,\ \binom{i+j}{j}=\sum_{k=0}^{i}\binom{i}{k}\binom{j}{k}$.
			\item En déduire que l'on peut écrire $B_n=A_nA_n^T$ où $A_n\in\mathcal{M}_{n+1}(\R)$ est triangulaire supérieure à coefficients diagonaux tous égaux à $1$.
			\item Conclure.
		\end{enumerate}
	\end{enumerate}
\end{ex}

\begin{ex}[]{}
	Soit $A$ et $B$ deux éléments de $\mathcal{M}_n(\K)$. Pour $j\in\llbracket1,n\rrbracket$, on note $A_j$ la matrice dont toutes les colonnes sont celles de $A$, sauf la première qui est la $j$-ème colonne de $B$, et $B_j$ la matrice dont toutes les colonnes sont celles de $B$, sauf la $j$-ème qui est la première colonne de $A$. Montrer la formule : $$\det(A)\det(B)=\sum_{j=1}^{n}\det(A_j)\det(B_j)$$
\end{ex}
\begin{proof}
	Fixer la matrice $B$ et montrer que le membre de droite est une forme $n$-linéaire alternée sur $\K^n$. Conclure en regardant la valeur sur la base canonique. On pourra reconnaître un développement par rapport à la première ligne (en développant les $\det(B_j)$ par rapport à la $j$-ème colonne).
\end{proof}

\section{Comatrice}

\begin{ex}[Moulin.Mat.77]{}
	Résoudre dans $\mathcal{M}_n(\R)$ : 
	\begin{enumerate}
		\item $A=\text{Com}(A)$ ;
		\item $A=(\text{Com}(A))^T$.
	\end{enumerate}
\end{ex}

\begin{ex}[Moulin.Mat.78]{}
	\begin{enumerate}
		\item $\Delta$ Déterminer le rang de $\text{Com}(A)^T$ en fonction du rang de $A$.
		\item Calculer $\det(\text{Com}(A)^T)$ en fonction de $\det(A)$.
		\item Exprimer $\text{Com}(AB)^T$ en fonction de $\text{Com}(A)^T$ et $\text{Com}(B)^T$ lorsque $A$ et $B$ sont inversibles. Que vaut $\text{Com}(A\inv)^T$ ?
		\item Calculer $\text{Com}(\text{Com}(A)^T)^T$.
	\end{enumerate}
\end{ex}

\begin{ex}[Moulin.Mat.79]{$\Delta\star$ Comatrice du produit}
	Soit $A, B \in \mathcal{M}_n(\C)$. Montrer que $$\text{Com}(AB)=\text{Com}(A)\text{Com}(B)$$
\end{ex}
\begin{proof}
	Commencer par le cas où les deux matrices sont inversibles. Le résultat s'obtient alors rapidement avec la formule des comatrices en fonction des déterminants et des transposées. Pour le cas général, utiliser la densité des matrices inversibles dans les matrices (définir les matrices $A_z=A+zI_n$ et $B_z=B+zI_n$) puis définir des polynômes coefficient par coefficient pour les matrices $\text{Com}(A_z B_z)$ et $\text{Com}(A_z)Com(B_z)$ et montrer qu'ils sont égaux car ils le sont en une infinité de points. Enfin, évaluer en 0 pour conclure.
\end{proof}

\section{TD Châteaux}
\begin{ex}{Formes linéaires sur $\mathcal{M}_n(\K)$}
	\begin{enumerate}
		\item Soit $\varphi$ une forme linéaire sur $\mathcal{M}_n(\K)$. Montrer qu'il existe une unique matrice $A\in\mathcal{M}_n(\K)$ telle que $\forall M\in\mathcal{M}_n(\K),\ \varphi(M)=\Tr(AM)$.
		\item Soit $\varphi$ une forme linéaire sur $\mathcal{M}_n(\K)$ telle que $\forall (A,B)\in\mathcal{M}_n(\K)^2,\ \varphi(AB)=\varphi(BA)$. montrer que $\varphi$ est proportionnelle à la trace.
		\item Soit $\varphi$ une forme linéaire sur $\mathcal{M}_n(\K)$ telle que $\forall(A,B,C)\in\mathcal{M}_n(\K)^3,\ \varphi(ABC)=\varphi(CBA)$. On suppose $n\geqslant 2$. Prouver que $\varphi=0$.
		\item On suppose que $\K$ est de caractéristique nulle. Soit $M\in\mathcal{M}_n(\K)$. Montrer l'équivalence : $$\Tr(M)=0\iff\exists(A,B)\in\mathcal{M}_n(\K),\ M=AB-BA$$ Pour cela : 
		\begin{enumerate}[label=\alph*.]
			\item Si $\Tr(M)=0$ et $M\ne 0$, montrer que $M$ est semblable à une matrice dont la première colonne est $(0\ 1\ 0\ \ldots\ 0)^T$.
			\item En déduire que si $\Tr(M)=0$, alors $M$ est semblable à une matrice dont la diagonale est nulle.
			\item Montrer l'implication attendue pour $M$ de diagonale nulle (en cherchant $A$ sous forme diagonale).
			\item Conclure.
		\end{enumerate}
		\item Démontrer que tout hyperplan de $\mathcal{M}_n(\K)$ contient des matrices inversibles. Indication : on pourra utiliser la première question.
	\end{enumerate}
\end{ex}
\begin{ex}{Transvections et dilatations}
	Soit $E$ un $\K$-espace vectoriel de dimension $n$. Soit $u$ un endomorphisme tel que $\text{rg}(u-\text{id})=1$.
	\begin{enumerate}
		\item Justifier l'existence d'un vecteur $a$ non nul et d'une forme linéaire $l$ non nulle tels que $\forall x\in E,\ u(x)=x+l(x)a$. Que vaut $\det(u)$ ?
		\item On suppose $l(a)=0$. Montrer que $\im(u-\text{id})\subset\ker(u-\text{id})$. On dit que $u$ est une transvection.
		\item On suppose $l(a)\ne 0$. Montrer que $\im(u-\text{id})$ et $\ker(u-\text{id})$ sont supplémentaires dans $E$. On dit alors que $u$ est une dilatation.
		\item On pose, pour $i$ et $j$ distincts compris entre $1$ et $n$, et pour $\lambda\in\K$, $T_{i,j}(\lambda)=I_n+\lambda E_{i,j}$. On note $\mathcal{T}$ l'ensemble des matrices de cette forme. On pose, pour $a\ne 0$, $D_a=I_n+(a-1)E_{n,n}$.
		\begin{enumerate}[label=\alph*.]
			\item Calculer $T_{i,j}(\lambda)T_{k,l}(\mu)$.
			\item Montrer que tout matrice inversible de déterminant $1$ est produit de matrices appartenant à $\mathcal{T}$.
			\item Montrer que toute matrice inversible est produit de matrices de $\mathcal{T}$ et d'une matrice $D_a$.
		\end{enumerate}
		\item Montrer que le groupe spécial linéaire $SL(E)$ est engendré par les transvections, puis que le groupe linéaire $GL(E)$ est engendré par les dilatations et les transvections.
	\end{enumerate}
\end{ex}
\begin{ex}{Une application des transvections}
	Soit $f:\mathcal{M}_n(\K)\to\K$ non constante telle que $\forall(A,B)\in\mathcal{M}_n(\K),\ f(AB)=f(A)f(B)$.
	\begin{enumerate}
		\item Montrer l'équivalence : $A\in\mathcal{GL}_n(\K)\iff f(A)\ne 0$.
		\item Désormais, et ce pour toute la suite de l'exercice, on suppose que $f$ est polynomiale par rapport aux coefficients de la matrice. Montrer que $f(T_{i,j}(\lambda))=1$ pour tous $i$ et $j$ distincts et tout $\lambda\in\K\st$.
		\item Montrer qu'il existe $r\in\N$ tel que $\forall a\in\K,\ f(D_a)=a^r$.
		\item En déduire que $\forall A\in\mathcal{M}_n(\K),\ f(A)=\det(A)^r$.
	\end{enumerate}
\end{ex}
\begin{ex}{Théorème de Burnside}
	On dit qu'un groupe $G$ est d'exposant fini lorsqu'il existe $d\in\N\st$ tel que pour tout $g\in G$, on ait $g^d=e$. Dans ce cas, on appelle exposant de $G$ le plus petit entier naturel non nul $d$ vérifiant cette propriété. On se propose ici de démontrer le théorème de Burnside : un sous-groupe de $\mathcal{GL}_n(\C)$ est fini si, et seulement s'il est d'exposant fini.
	\begin{enumerate}
		\item Soit $A\in\mathcal{M}_n(\C)$ telle que $\forall k\in\N,\ \Tr(A^k)=0$. montrer que $A$ est nilpotente.
		\item Soit $G$ un sous-groupe de $\mathcal{GL}_n(\C)$. Soit $(M_i)_{1\leqslant i\leqslant m}$ une base de $\text{Vect}(G)$ constituée d'éléments de $G$. On définit $f$ de $G$ dans $\C^m$ par : $$f(M)=\left(\Tr(MM_i)\right)_{1\leqslant i\leqslant m}$$ Démontrer que si $f(A)=f(B)$, alors $AB\inv-I_n$ est nilpotente. 
		\item En déduire que $f$ est injective lorsque $G$ est d'exposant fini.
		\item Conclure.
	\end{enumerate}
\end{ex}
\begin{proof}
	\begin{enumerate}
		\item Classique : un déterminant de Vandermonde fait l'affaire.
		\item Remarquons que $(M_i)$ existe bel et bien par le théorème de la base extraite. Supposons $f(A)=f(B)$. Alors : $$\forall i\in\llbracket1,m\rrbracket,\ \Tr((A-B)M_i)=0$$ On en déduit $$\forall M\in \text{Vect}(G),\ \Tr(AM)=\Tr(BM)$$ On pose alors $U=AB\inv-I_n$. On procède par récurrence en remarquant que $$U^{k+1}=AB\inv U^k-U^k$$ avec $B\inv U^k\in\text{Vect}(G)$ donc on en déduit que $\Tr(U^k)=0$ pour tout $k$ et donc que $U$ est nilpotente.
		\item Il suffit de montrer que $U$ est diagonalisable, et donc il suffit de montrer que $AB\inv$ l'est. Supposons $G$ d'exposant fini. Comme $AB\inv$ est annulé par $X^d-1$ (scindé à racines simples sur $\C$) où $d$ est l'exposant de $G$, on en déduit que $AB\inv$ est diagonalisable puis que $U=0$ et donc que $f$ est injective.
		\item Supposons toujours $G$ d'exposant fini (le sens direct est immédiat). Comme les éléments de $G$ sont tous annulés par $X^d-1$, leurs valeurs propres sont toutes des éléments de $\U_d$ donc, comme les $MM_i$ sont des éléments de $G$, $f$ prend un nombre fini de valeurs (la trace est la somme des valeurs propres). Comme $f$ est injective, on en déduit que $G$ est fini.
	\end{enumerate}
\end{proof}

\end{document}